\section{第2章\quad 随机向量随机数的抽样方法}

\subsection{故事：从一维到多维，难在"相关"}

第 1 章我们学会了造一维随机数。多维呢？如果各分量\textbf{独立}，那太简单了——
每个分量单独抽就行。真正的麻烦在于分量之间\textbf{有相关性}：
比如身高和体重是正相关的，不能各抽各的。
本章三种思路对付相关性：
\begin{itemize}[leftmargin=2em]
  \item \textbf{变换法}：从独立分量出发，做一个线性（或非线性）变换"注入"相关性；
  \item \textbf{条件分布法}：一个一个分量地抽，后抽的"看着"先抽的来（链式法则）；
  \item \textbf{舍选法}：一维舍选的多维翻版。
\end{itemize}

\subsection{变换抽样法}

\paragraph{基础铺垫：多维密度变换公式。}
设 $p$ 维随机向量 $\bm\xi$ 有密度 $f(x_1,\dots,x_p)$，作可逆变换 $\bm\eta=g(\bm\xi)$，则 $\bm\eta$ 的密度为
\[
p(\bm y)=f\big(g^{-1}(\bm y)\big)\cdot|J| ,\qquad
J=\det\Big(\frac{\partial x_i}{\partial y_j}\Big),
\]
即"原密度 $\times$ 逆变换 Jacobi 行列式的绝对值"。
直觉：变换会拉伸/压缩空间，$|J|$ 是体积缩放因子，密度要按它折算才能保证总概率为 1。
（一维的 $p_Y(y)=p_X(x)\big|\frac{\d x}{\d y}\big|$ 就是它的特例；Box--Muller 的正确性也由它证明。）

\paragraph{主角：多维正态。}
多维正态 $N_p(\bm\mu,\Sigma)$ 是统计中最重要的多维分布，它有个极好的性质：
\textbf{独立标准正态的线性变换仍是正态}，且均值协方差变换规律简单：
\[
\bm X=\bm\mu+L\bm Z\ \Longrightarrow\ \bm X\sim N_p(\bm\mu,\ LL^{T}).
\]
所以只要把 $\Sigma$ 分解成 $LL^T$（\textbf{Cholesky 分解}，$L$ 为下三角矩阵），
再抽独立标准正态向量 $\bm Z$，线性组合一下就完工。

\begin{pbox}
类比一维：一维里我们用 $X=\mu+\sigma Z$ 把标准正态"平移+拉伸"成任意正态；
多维里 $L$ 就是 $\sigma$ 的矩阵版——Cholesky 分解 $\Sigma=LL^T$ 相当于给协方差矩阵"开平方"。
二维考试只需要背一个结果：$\eta_2=\mu_2+\sigma_2(\rho\,\xi_1+\sqrt{1-\rho^2}\,\xi_2)$。
看结构：$\rho$ 那份是"拄着 $\xi_1$ 的部分"（制造相关），$\sqrt{1-\rho^2}$ 那份是自己的
"新鲜随机性"，两者平方和为 1，保证总方差不变。
\end{pbox}

\begin{ebox}{生成 $N_p(\bm\mu,\Sigma)$ 的变换法}
\begin{enumerate}[leftmargin=2em]
  \item 对 $\Sigma$ 做 Cholesky 分解 $\Sigma=LL^T$；
  \item 用 Box--Muller 生成独立 $Z_1,\dots,Z_p\sim N(0,1)$；
  \item 输出 $\bm X=\bm\mu+L\bm Z$。
\end{enumerate}
正确性：$\Cov(\bm X)=L\,\Cov(\bm Z)\,L^T=L I L^T=\Sigma$。
\end{ebox}

\begin{exbox}[二维正态的变换法（课件原题）]
生成 $N_2(\bm\mu,\Sigma)$，其中 $\mu_1=\mu_2=0$，$\sigma_1=\sigma_2=2$，$\rho=0.8$。
写出变换公式并验证。
\end{exbox}
\begin{sol}
二维情形的"手工 Cholesky"是必背公式：取独立 $\xi_1,\xi_2\sim N(0,1)$，令
\[
\eta_1=\mu_1+\sigma_1\xi_1,\qquad
\eta_2=\mu_2+\sigma_2\big(\rho\,\xi_1+\sqrt{1-\rho^2}\,\xi_2\big).
\]
代入数值：
\[
\eta_1=2\xi_1,\qquad
\eta_2=2(0.8\xi_1+\sqrt{1-0.64}\,\xi_2)=1.6\xi_1+1.2\xi_2 .
\]
\textbf{验证}：
$\Var\eta_1=4$；$\Var\eta_2=1.6^2+1.2^2=2.56+1.44=4$；
$\Cov(\eta_1,\eta_2)=2\times1.6\,\Var\xi_1=3.2=\sigma_1\sigma_2\rho=2\cdot2\cdot0.8$。
均值、方差、相关系数全部正确，且正态线性组合仍正态，故 $(\eta_1,\eta_2)\sim N_2(\bm\mu,\Sigma)$。
\end{sol}

\subsection{条件分布法：链式法则的胜利}

\paragraph{直觉先行。}
联合密度永远可以"剥洋葱"：
\[
f(x_1,x_2,\dots,x_p)=f_1(x_1)\cdot f_2(x_2\mid x_1)\cdots f_p(x_p\mid x_1,\dots,x_{p-1}).
\]
所以抽一个 $p$ 维向量 = 抽 $p$ 次一维：先抽 $x_1$，
然后"在 $x_1$ 已知的世界里"抽 $x_2$，依此类推。
相关性自动被条件分布编码进去了。

\begin{ebox}{算法步骤}
\begin{enumerate}[leftmargin=2em]
  \item 从边际 $f_1(x_1)$ 抽 $x_1$；
  \item 从条件密度 $f_2(x_2\mid x_1)$ 抽 $x_2$；
  \item \dots 直到 $f_p(x_p\mid x_1,\dots,x_{p-1})$ 抽出 $x_p$，输出整个向量。
\end{enumerate}
\end{ebox}

\paragraph{多维正态的条件分布公式（需记）。}
把 $\bm X,\bm\mu,\Sigma$ 分块：
$\bm X=\binom{\bm X_1}{\bm X_2}$，
$\Sigma=\begin{pmatrix}\Sigma_{11}&\Sigma_{12}\\\Sigma_{21}&\Sigma_{22}\end{pmatrix}$，则
\[
\bm X_2\mid\bm X_1=\bm x_1\ \sim\ N\big(\bm\mu_2+\Sigma_{21}\Sigma_{11}^{-1}(\bm x_1-\bm\mu_1),\ \ \Sigma_{22}-\Sigma_{21}\Sigma_{11}^{-1}\Sigma_{12}\big).
\]

\begin{pbox}
用"身高体重"记二维版：已知一个人身高 $x_1$，猜他体重。
猜测中心不再是全人群平均 $\mu_2$，而要按身高偏高/偏矮的程度修正：
$\mu_2+\rho\frac{\sigma_2}{\sigma_1}(x_1-\mu_1)$（身高越高、相关越强，修正越多）；
同时因为知道了身高，不确定性变小：方差从 $\sigma_2^2$ 缩到 $\sigma_2^2(1-\rho^2)$。
记住这两句话，公式就背下来了。
\end{pbox}

\begin{exbox}[条件分布法生成二维正态]
用条件分布法生成 $N_2(\bm\mu,\Sigma)$（参数同上例：$\sigma_1=\sigma_2=2,\rho=0.8,\bm\mu=\bm0$）。
\end{exbox}
\begin{sol}
二维情形条件公式化简为：
\[
X_1\sim N(\mu_1,\sigma_1^2),\qquad
X_2\mid X_1=x_1\sim N\Big(\mu_2+\rho\frac{\sigma_2}{\sigma_1}(x_1-\mu_1),\ \sigma_2^2(1-\rho^2)\Big).
\]
代入数值：$X_1\sim N(0,4)$；$X_2\mid X_1=x_1\sim N(0.8x_1,\ 4\times0.36)=N(0.8x_1,\,1.44)$。

\textbf{步骤}：(i) 抽 $z_1\sim N(0,1)$，令 $x_1=2z_1$；
(ii) 抽 $z_2\sim N(0,1)$，令 $x_2=0.8x_1+1.2z_2$。
展开发现 $x_2=1.6z_1+1.2z_2$——与变换法完全一致！两种方法殊途同归。
\end{sol}

\subsection{多维舍选法}

与一维完全平行：目标密度 $p(\bm x)$，建议密度 $f(\bm x)$（常取各分量独立、容易抽的分布），
$M\ge\sup p/f$；抽 $\bm y\sim f$、$u\sim U(0,1)$，$u\le\frac{p(\bm y)}{Mf(\bm y)}$ 则接受。

\begin{exbox}[单位圆盘上的均匀点]
设计算法在单位圆盘 $\{x^2+y^2\le1\}$ 上抽均匀分布的点，并求接受概率。
\end{exbox}
\begin{sol}
在正方形 $[-1,1]^2$ 上抽均匀点很容易：$x=2u_1-1,\ y=2u_2-1$。
圆盘含于正方形内，于是：
\textbf{(i)} 产生 $u_1,u_2\sim U(0,1)$，令 $x=2u_1-1,y=2u_2-1$；
\textbf{(ii)} 若 $x^2+y^2\le1$ 接受 $(x,y)$，否则返回 (i)。

接受概率 $=$ 圆盘面积/正方形面积 $=\pi/4\approx78.5\%$。
被接受的点在圆盘内自动均匀（均匀点限制在子区域上仍均匀）。
\end{sol}

\subsection{离散随机向量：条件分布法大显身手}

\begin{exbox}[放球问题（课件原题）]
把 $n$ 个球独立地等概率放入 A、B、C 三个篮子。
记 $X$ 为落入 A 的个数、$Y$ 为落入 B 的个数。给出生成 $(X,Y)$ 的方法。
\end{exbox}
\begin{sol}
边际：每个球进 A 的概率 $1/3$，故 $X\sim B(n,\frac13)$。
条件：已知 $X=x$，剩下 $n-x$ 个球只在 B、C 中等可能选择，进 B 概率
$\frac{1/3}{2/3}=\frac12$，故 $Y\mid X=x\sim B(n-x,\frac12)$。

\textbf{步骤}：(i) 抽 $x\sim B(n,1/3)$；(ii) 抽 $y\sim B(n-x,1/2)$；输出 $(x,y)$。
\end{sol}

\paragraph{推广：多项分布 $\mathrm{Mult}(n;p_1,\dots,p_r)$。}
有放回取 $n$ 次、每次落入第 $i$ 类概率 $p_i$，$X_i$ 为第 $i$ 类计数。两种生成方式：
\begin{enumerate}[leftmargin=2em]
  \item \textbf{逐次分类法}：做 $n$ 次离散查表抽样（第 1 章方法），统计各类频数；
  \item \textbf{条件分布法}：
  \[
  X_1\sim B(n,p_1),\qquad
  X_k\mid X_1,\dots,X_{k-1}\sim B\Big(n-\sum_{i<k}x_i,\ \frac{p_k}{1-\sum_{i<k}p_i}\Big).
  \]
  直觉：抽掉前几类后，剩余的球在剩余类别中按"重新归一化"的概率分配。
\end{enumerate}

\subsection{本章考点地图}

\begin{ebox}{高频题型清单}
\begin{enumerate}[leftmargin=2em]
  \item \textbf{变换法生成二维正态}：默写 $\eta_1,\eta_2$ 公式，代数值并验证三个矩（如 $\sigma_1=\sigma_2=2,\rho=0.8$）；
  \item \textbf{条件分布法生成二维正态}：默写条件正态参数 $\mu_2+\rho\frac{\sigma_2}{\sigma_1}(x_1-\mu_1)$、$\sigma_2^2(1-\rho^2)$，写两步算法；
  \item \textbf{离散向量条件分布法}：放球/多项分布的逐分量条件二项分布；
  \item \textbf{多维舍选}：不规则区域上的均匀抽样+接受概率（面积比）；
  \item \textbf{Jacobi 公式}：用密度变换公式验证某变换的正确性。
\end{enumerate}
\end{ebox}
